\courseheader{MK1}{Organizational Behaviour \& Managing People}
\begin{minipage}[t]{.66\textwidth}\vspace{0pt}
\textcolor{navy}{\large\bfseries Learning Outcomes}\par
\begin{itemize}
\item Understand how to collaborate with other parties.
\item Concept of conflict management to solve problems in a small group learning environment.
\item Demonstrate constructive feedback in a small group learning environment.
\item Demonstrate the tendency to take the initiative.
\item Students are able to influence others in their team.
\item Inspire and empower others by evaluating, analyzing, and providing criticism about how leadership behavior and characteristics affect workers and achievements in business.
\item Demonstrate the ability to manage change.
\end{itemize}
\end{minipage}\hfill
\begin{minipage}[t]{.32\textwidth}\vspace{0pt}
\staffbox{\textbf{Muhammad Yorga Permana, Ph.D.}\par Widhyawan Prawiraatmadja, Ph.D.\par Dr. Ir. Madju Yuni Ros Bangun\par Yudo Anggoro, Ph.D.}{Efryta Wulan (081245103149)}
\end{minipage}
\rowcolors{2}{paleblue}{white}
\begin{tabularx}{\textwidth}{M{1.8cm}YY}\tableheader
1&Understanding Self \& Leading Others&Strategizing Your Life: Self-Awareness, Values \& Purpose\\
2&Understanding Self \& Leading Others&Who Am I at Work? Identity, Personality \& Individual Differences\\
3&Understanding Self \& Leading Others&Motivation \& Human Behavior at Work\\
4&Understanding Self \& Leading Others&Leadership \& Influence\\
5&Navigating Teams \& Organizations&Team Dynamics, Communication, Negotiation \& Conflict\\
6&Navigating Teams \& Organizations&Organizational Culture\\
7&Navigating Teams \& Organizations&Power \& Politics in Organizations\\
8&Building the Organization \& Talent&People \& Strategy: Human Capital as Competitive Advantage\\
9&Building the Organization \& Talent&Organization Design \& Strategic Workforce Planning\\
10&Building the Organization \& Talent&Strategic Talent Acquisition \& Employer Branding\\
11&Building the Organization \& Talent&Talent Management, Development \& Succession\\
12&Driving Performance \& Transformation&Performance Management, Compensation, and Benefit\\
13&Driving Performance \& Transformation&Leading Change, Transformation \& the Future of Workplace\\
14&Driving Performance \& Transformation&Employee Experience, Engagement \& Inclusion\\
\end{tabularx}\rowcolors{2}{}{}
\begin{visualexplanation}
The curriculum moves from the individual to the organization and then to transformation. Weeks 1--4 build self-awareness, motivation, leadership, and influence; Weeks 5--7 address team and organizational dynamics; Weeks 8--11 build organization and talent systems; Weeks 12--14 focus on performance, change, and employee experience. This sequence directly supports the listed outcomes concerning collaboration, conflict management, feedback, initiative, influence, empowerment, and change management. The teaching team and course assistant shown above the table support the full 14-week progression.
\end{visualexplanation}

\courseheader{MK2}{Financial Management and Strategy}
\begin{minipage}[t]{.66\textwidth}\vspace{0pt}
\textcolor{navy}{\large\bfseries Learning Outcomes}\par
\begin{itemize}
\item Demonstrate understanding of the financial management function and the firm's financial goals.
\item Conduct analytical reviews of financial statements and projections; apply time-value-of-money concepts and techniques.
\item Apply tools to analyze financial decisions---investment, financing, payout, and working capital---and address real-world financial issues.
\item Identify the financial environment and institutions; recognize the importance of ethics in financial decision making.
\end{itemize}
\end{minipage}\hfill
\begin{minipage}[t]{.32\textwidth}\vspace{0pt}
\staffbox{\textbf{Dr. Arif Zamani}\par Dr. Subiakto Soekarno\par Dr. Aswin Rahadi\par Yunieta Anny Nainggolan, Ph.D.\par Oktofa Yudha Sudrajat, Ph.D.}{Annisa Rizkia Syaputri (081809905646)}
\end{minipage}
\rowcolors{2}{paleblue}{white}
\begin{tabularx}{\textwidth}{M{1.8cm}YY}\tableheader
1&Financial Foundations&Financial Management Function \& Objectives\\
2&Financial Foundations&Financial Statement Analysis\\
3&Financial Foundations&Financial Statement Projections\\
4&Cost \& Managerial Accounting&Management Accounting Fundamentals\\
5&Cost \& Managerial Accounting&Cost Analysis \& Activity-Based Costing\\
6&Time Value of Money&Time Value of Money (1)\\
7&Time Value of Money&Time Value of Money (2)\\
8&Corporate Financial Decisions&Capital Budgeting (1)\\
9&Corporate Financial Decisions&Capital Budgeting (2)\\
10&Corporate Financial Decisions&Pay-Out Policy\\
11&Corporate Financial Decisions&Financing Decisions \& Financial Environment\\
12&Corporate Financial Decisions&Working Capital Management\\
13&Financial Risk Management&Definition of risk, risk-return relationship, \& risk governance\\
14&Ethics \& Synthesis&Group Project Presentations\\
\end{tabularx}\rowcolors{2}{}{}
\begin{visualexplanation}
The 14-week sequence begins with financial foundations and statement analysis, then introduces managerial accounting and time value of money before moving into capital budgeting, payout, financing, and working-capital decisions. Risk governance in Week 13 and group presentations in Week 14 complete the progression with ethics and synthesis. The table therefore maps directly to the outcomes: understanding the finance function, reviewing statements and projections, applying decision tools, recognizing the financial environment, and incorporating ethics into decisions.
\end{visualexplanation}

\courseheader{MK3}{Marketing Management}
\begin{minipage}[t]{.66\textwidth}\vspace{0pt}
\textcolor{navy}{\large\bfseries Learning Outcomes}\par
\begin{itemize}
\item Identifying and analyzing the core marketing problems in a business case.
\item Identifying ethical marketing issues in a business case.
\item Taking initiative in the process of designing a Marketing Plan.
\item In groups, students are able to identify and analyze the core marketing problems in the process of evaluating a Marketing Plan design.
\end{itemize}
\end{minipage}\hfill
\begin{minipage}[t]{.32\textwidth}\vspace{0pt}
\staffbox{\textbf{Ira Fachira, Ph.D.}\par Atik Aprianingsih, DBA\par Ilma Aulia Zaim, Ph.D.\par Nila Armelia, Ph.D.\par Dr. Gallang Perdhana Dalimunthe}{Nia Desiana (081214533663)}
\end{minipage}
\rowcolors{2}{paleblue}{white}
\begin{tabularx}{\textwidth}{M{1.8cm}YY}\tableheader
1&Foundations&Creating Customer Value\\
2&Foundations&Analyzing the Marketing Environment\\
3&Customer \& Market Insight&Managing Marketing Information \& Customer Insights (1)\\
4&Customer \& Market Insight&Managing Marketing Information \& Customer Insights (2)\\
5&Customer \& Market Insight&Marketing Research\\
6&STP Strategy&Segmentation \& Targeting\\
7&STP Strategy&Positioning\\
8&STP Strategy&Games: Selling\\
9&Marketing Mix (4Ps)&Marketing Mix --- Product \& Services\\
10&Marketing Mix (4Ps)&Marketing Mix --- Price\\
11&Marketing Mix (4Ps)&Marketing Mix --- Place (Channel) \& Digital Marketing\\
12&Marketing Mix (4Ps)&Marketing Mix --- Promotion\\
13&Branding \& Synthesis&Branding\\
14&Branding \& Synthesis&Special Topic / Guest Lecture\\
\end{tabularx}\rowcolors{2}{}{}
\begin{visualexplanation}
The curriculum starts with customer value and the marketing environment, develops customer and market insight through information and research, and then uses segmentation, targeting, positioning, and a selling game to frame strategic choices. Weeks 9--12 cover the four marketing-mix levers, followed by branding and a special topic or guest lecture. This flow supplies the knowledge required to identify core and ethical marketing problems, take initiative in designing a Marketing Plan, and evaluate that design collaboratively.
\end{visualexplanation}

\courseheader{MK4}{Operations and Supply Chain Management}
\begin{minipage}[t]{.66\textwidth}\vspace{0pt}
\textcolor{navy}{\large\bfseries Learning Outcomes}\par
\begin{itemize}
\item Students will be able to demonstrate an understanding of theory, concepts, models, or reasoning in operations and supply chain management related to specific problems/issues/contexts as discussed throughout the course.
\item Students will be able to critique concepts, models, or reasoning in operations and supply chain management, so as to suggest solutions and managerial recommendations.
\item Students will understand the role of operations and supply chain management within the global business context.
\end{itemize}
\end{minipage}\hfill
\begin{minipage}[t]{.32\textwidth}\vspace{0pt}
\staffbox{\textbf{Noorhan Firdaus Pambudi, Ph.D.}\par Dr. Yuanita Handayati\par Berty Argiyantari, Ph.D.\par Dr. Liane Okdinawati}{Laisa Aritonang (081221712180)}
\end{minipage}
\rowcolors{2}{paleblue}{white}
\begin{tabularx}{\textwidth}{M{1.8cm}YY}\tableheader
1&Strategy and Product Design&Introduction to Operations and Supply Chain Management\\
2&Strategy and Product Design&Strategy\\
3&Strategy and Product Design&Product and Service Design\\
4&Strategy and Product Design&Projects\\
5&Capacity and Process&Strategic Capacity Planning\\
6&Capacity and Process&Process Design and Analysis\\
7&Capacity and Process&Waiting Line Analysis and Simulation\\
8&Quality Assurance&Lean Operations\\
9&Quality Assurance&Lean Six Sigma\\
10&Quality Assurance&Statistical Quality Control\\
11&Demand Fulfillment&Forecasting\\
12&Demand Fulfillment&Sales \& Operations Planning\\
13&Demand Fulfillment&Inventory Management\\
14&Demand Fulfillment&MRP and ERP\\
\end{tabularx}\rowcolors{2}{}{}
\begin{visualexplanation}
The table groups the course into four operational layers. Weeks 1--4 align strategy with product, service, and project design. Weeks 5--7 address capacity and process behavior. Weeks 8--10 provide three approaches to quality assurance: lean operations, Lean Six Sigma, and statistical quality control. Weeks 11--14 connect forecasting, sales and operations planning, inventory, MRP, and ERP to demand fulfillment. Together these blocks let students understand and critique operations and supply-chain models, formulate managerial recommendations, and place the function in a global business context.
\end{visualexplanation}

\courseheader{MK5}{Problem Solving, Decision Making and Negotiation}
\begin{minipage}[t]{.66\textwidth}\vspace{0pt}
\textcolor{navy}{\large\bfseries Learning Outcomes}\par
\begin{itemize}
\item Identifying and analyzing the core problems in a business case.
\item Building multiple perspectives by integrating various factors, relevant business functions, \& contextual information.
\item Developing solution options based on the perspectives built and the intended vision.
\item Making decisions based on the solution options developed.
\item Collaborating effectively with others.
\item Applying conflict management concepts to resolve problems within a team.
\item Providing constructive feedback within a team.
\item Communicating persuasively to influence and convince others.
\end{itemize}
\end{minipage}\hfill
\begin{minipage}[t]{.32\textwidth}\vspace{0pt}
\staffbox{\textbf{Dr. Valid Hasyimi}\par Prof. Dr. Ir. Utomo Sarjono Putro\par Prof. Dr.Eng. Pri Hermawan\par Santi Novani, Ph.D.}{Syifa El-Baehaqi (08999230348)}
\end{minipage}
\rowcolors{2}{paleblue}{white}
\begin{tabularx}{\textwidth}{M{1.8cm}YY}\tableheader
1&Foundations&Introduction to Negotiation\\
2&Foundations&Rational Framework of Decision Making\\
3&Kepner-Tregoe&Kepner-Tregoe Part 1\\
4&Kepner-Tregoe&Kepner-Tregoe Part 2\\
5&Kepner-Tregoe&Decision Analysis: SMART\\
6&Kepner-Tregoe&Decision Analysis: AHP and Decision Tree\\
7&Scenario Planning and Systems Thinking&Scenario Planning and Systems Thinking (1)\\
8&Scenario Planning and Systems Thinking&Scenario Planning and Systems Thinking (2)\\
9&Distributive Negotiation&Negotiation Dilemma\\
10&Distributive Negotiation&Negotiation Style\\
11&Distributive Negotiation&Claiming value (``Slicing the Pie'') in negotiation\\
12&Integrative Negotiation&From distributive to integrative negotiation\\
13&Integrative Negotiation&Integrative Negotiation Strategy and Technique\\
14&Integrative Negotiation&Multiparty Negotiation \& Wrap Up\\
\end{tabularx}\rowcolors{2}{}{}
\begin{visualexplanation}
The course first establishes negotiation and rational decision-making foundations, then devotes four weeks to Kepner-Tregoe methods, including SMART, AHP, and decision trees. Two weeks of scenario planning and systems thinking expand the range of perspectives. The final six weeks move from distributive negotiation dilemmas and value claiming to integrative, multiparty negotiation. That progression connects problem identification and option development to decisions, collaboration, conflict management, constructive feedback, and persuasive communication.
\end{visualexplanation}

\courseheader{MK6}{Applied Business Analytics}
\begin{minipage}[t]{.66\textwidth}\vspace{0pt}
\textcolor{navy}{\large\bfseries Learning Outcomes}\par
\begin{itemize}
\item Identifying and analyzing the core problems in a business case.
\item Building multiple perspectives by integrating various factors, related business functions, and relevant contextual information.
\item Developing solution options based on the perspectives built and the intended vision.
\item Making decisions based on the solution options developed.
\end{itemize}
\end{minipage}\hfill
\begin{minipage}[t]{.32\textwidth}\vspace{0pt}
\staffbox{\textbf{Dr. Eng. Manahan Siallagan}\par Dr. Shimaditya Nuraeni\par M. Apriandito Arya Saputra, MBA\par Bagus Rullym MM}{Tyas Larassati (081809230368)}
\end{minipage}
\rowcolors{2}{paleblue}{white}
\begin{tabularx}{\textwidth}{M{1.8cm}YY}\tableheader
1&Foundations&Introduction to Business Analytics\\
2&Foundations&Data-Driven Decision Making\\
3&Foundations&Predictive Analytics for Business\\
4&Applied Analytics&Customer Analytics\\
5&Applied Analytics&Marketing Analytics\\
6&Applied Analytics&Financial Analytics\\
7&Applied Analytics&Supply Chain Analytics\\
8&Applied Analytics&Risk Analytics\\
9&Data \& Responsibility&Big Data in Business\\
10&Data \& Responsibility&Ethical Considerations in Business Analytics\\
11&Data \& Responsibility&Practical: Big Data Technologies\\
12&Data \& Responsibility&Practical: Data Privacy and Security\\
13&Advanced Analytics&Business Intelligence \& Visualization\\
14&Advanced Analytics&Future Trends in Business Analytics\\
\end{tabularx}\rowcolors{2}{}{}
\begin{visualexplanation}
The curriculum begins with analytics foundations, data-driven decisions, and prediction. Weeks 4--8 apply analytics to customers, marketing, finance, supply chains, and risk. Weeks 9--12 place those applications within big data, ethics, enabling technologies, privacy, and security. Business intelligence, visualization, and future trends close the course. The movement from foundations through domain applications and responsible data practice supports the outcomes of identifying problems, integrating business perspectives and context, developing solution options, and making decisions from those options.
\end{visualexplanation}